% Body, short version (IEEE, 8 pages including references)
%======================================================================
\section{Introduction}
%======================================================================

Typed decision models return an answer in a fixed format, such as a
label, a yes/no or a rating, together with a probability or a
confidence. The answer can be well formed and still wrong. Chain several
of these answers into a workflow and a mistake at one step becomes the
input to the next. Practitioners often assume that the confidence each
step reports can be combined into one workflow-level signal that says
when a person should look at a case.

Jev, TypeSafe's typed decision model, has been benchmarked on single
tasks \citep{deusser2026jev,meng2026jevaudit,denney2026jev}, and
uncertainty along multi-step agent trajectories has been studied
separately \citep{lin2026uqagents}. Our question is operational. When a
workflow chains several typed decisions, which signals identify the
cases where it fails, and do those signals find the costly errors that
justify a human review? The two parts can come apart. A signal can rank
failures well and still be of little use for choosing what to review.

Jev's primitives also report uncertainty in different forms. \choice{}
returns a probability for each label plus a scalar confidence, \score{}
returns an ordinal distribution plus a scalar confidence, and \noul{}
returns a single probability with no confidence field. TypeSafe
describes its confidence as one reasonable summary of a distribution and
advises setting thresholds on one's own data
\citep{typesafe2026confidence}. So we do not test a calibration claim.
We measure how these numbers behave as escalation signals once the
primitives are chained.

We ask three questions. RQ1: which signals rank workflow failures, the
model's confidence or what the workflow decided? RQ2: do those signals
stay useful when the label distribution shifts between development and
deployment? RQ3: do signals that rank failures also find the reviews that
reduce cost?

The contribution is methodological. We measure uncertainty at the level
of a workflow, where failure is defined on the action taken rather than
on any single answer. We separate ranking failures from reducing cost by
scoring a fixed review budget against a frozen cost matrix and an oracle
bound. And we pre-registered every confirmatory test, froze the
estimators by hash before their data existed, and replicated the main
result across five models and a second, chronologically split workflow.

In short, the workflow's decisions predict its failures much better than
its confidence does, and the gap holds across models and workflows. Most
of that signal is the error rate of the chosen action. Those error rates
shift with the label mix. And none of the tested signals made review pay
under our costs.

%======================================================================
\section{Related work}
\label{sec:related}
%======================================================================

\para{Jev and typed decisions} Deu{\ss}er et al.\ benchmarked Jev
zero-shot on 37 datasets and found well-calibrated \choice{}
probabilities, while binary probabilities ranked well but sat poorly
against a 0.5 threshold \citep{deusser2026jev}. A calibration audit found
that calibration changed with the question form and that one fitted
temperature removed most miscalibration \citep{meng2026jevaudit}, and a
study of political-science replications found Jev's calibration
advantage inconsistent \citep{denney2026jev}. Open instruction-tuned
models can produce Jev-like decisions from answer-label logits
\citep{geuter2026innerjev}, which is the readout we use for the general
models. All of this work evaluates single requests.

\para{Calibration, selection and deferral} Post-hoc calibration
\citep{platt1999probabilistic,guo2017calibration}, selective prediction
\citep{elyaniv2010selective,geifman2017selective} and learning to defer
\citep{madras2018predict,mozannar2020defer} all concern one decision.
Bansal et al.\ show that the most accurate model need not make the best
human--AI team \citep{bansal2021teammate}, and Park et al.\ argue that
error risk and review value are different quantities under a review
budget \citep{park2026riskreview}. Our cost result is a case of the
second point in a dependent workflow.

\para{Label shift and composition} Posteriors can be adjusted to new
class priors \citep{saerens2002adjusting}, and the priors can be
estimated from a black-box predictor's confusion matrix (BBSE)
\citep{lipton2018bbse}. Product rules for combining classifiers assume
independent, comparable scores \citep{kittler1998combining}. Conformal
methods have been composed over pipeline stages \citep{kotte2026pasc},
and step-level calibration does not guarantee trajectory-level
calibration for agents \citep{lin2026uqagents}. LLM raters compress
ordinal scales toward the middle \citep{zhang2026centraltendency}.

\para{Cascades and software reliability} Cascades route a single query
through model tiers and decide when to defer
\citep{zellinger2025early,bouchard2026escalation,jung2025trust}.
Component-based reliability models treat error propagation between
components as an architectural property that decides where detection
and recovery belong \citep{cortellessa2007errorprop}, and bug severity
prediction from report text is a long-studied triage task
\citep{lamkanfi2010severity}. Our workflow is a small instance of the
same propagation problem, with human escalation as the recovery step.

We found no prior work that tests, on real labels and held-out data,
which signals predict failure in a chain of structurally different typed
decisions and whether they find the reviews that pay.

%======================================================================
\section{Methodology}
\label{sec:setting}
%======================================================================

\begin{widefigure}
\centering
\resizebox{\linewidth}{!}{% Study architecture (shared by all wrappers): data -> typed decision workflow -> evaluation.
\begin{tikzpicture}[
  font=\footnotesize,
  src/.style={draw=black!45, fill=black!4, rounded corners=2pt, align=center, text width=2.7cm, inner sep=2pt, line width=0.5pt},
  prim/.style={draw=blue!60!black, fill=white, rounded corners=2pt, align=center, minimum width=2.05cm, minimum height=0.95cm, inner sep=2pt, line width=0.7pt, drop shadow={opacity=0.18, shadow xshift=0.4mm, shadow yshift=-0.4mm}},
  box/.style={draw=#1!65!black, fill=#1!9, rounded corners=2pt, align=center, inner sep=3pt, line width=0.5pt},
  arr/.style={-{Stealth[length=1.7mm]}, line width=0.65pt, draw=black!75},
  refarr/.style={-{Stealth[length=1.7mm]}, line width=0.55pt, draw=black!50, densely dashed},
  lane/.style={draw=#1!45, fill=#1!3, rounded corners=4pt, inner sep=2.2mm, line width=0.6pt},
  step/.style={circle, fill=#1!65!black, text=white, font=\scriptsize\bfseries, inner sep=1.1pt, minimum size=3.6mm},
  steplabel/.style={font=\footnotesize\bfseries, text=#1!60!black, anchor=west}]

% ---------------- (1) data ----------------
\node[src] (ker) {Linux kernel Bugzilla\\[-1pt]{\scriptsize 240 dev\,/\,250 test bugs}};
\node[src, above=2.5mm of ker] (moz) {Mozilla Bugzilla\\[-1pt]{\scriptsize 202 dev\,/\,250 test bugs}};
\node[src, below=2.5mm of ker] (k8s) {Kubernetes GitHub\\[-1pt]{\scriptsize 500 dev\,/\,582 test issues}};
\node[font=\scriptsize, text=black!65, text width=2.7cm, align=center, below=1.2mm of k8s] (cleannote) {all texts cleaned: privacy scrub, label-leak removal, $\leq$2{,}000 characters};
\foreach \t in {moz,k8s} \draw[line width=0.65pt, draw=black!75] (\t.east) -- ++(3mm,0) |- (ker.east);
\coordinate (bus) at ($(ker.east)+(3mm,0)$);

% ---------------- (2) workflow ----------------
\node[prim, right=15mm of ker] (c) {\choice{}\\[-1pt]routing $c_1$};
\node[prim, right=8.5mm of c] (n) {\noul{}\\[-1pt]bug\,/\,regression $c_2$};
\node[prim, right=8.5mm of n] (s) {\score{}\\[-1pt]severity\,/\,priority $c_3$};
\foreach \x in {c,n,s} \node[font=\scriptsize, text=blue!55!black, above=0.6mm of \x] {$3\times$ call $\rightarrow$ modal vote};
\draw[arr] (c) -- node[above, font=\scriptsize]{$x{+}c_1$} (n);
\draw[arr] (n) -- node[above, font=\scriptsize]{$x{+}c_{1:2}$} (s);
\path (n.south) -- (s.south) coordinate[pos=0.5] (ns);
\node[box=orange, below=9mm of ns, text width=5.2cm] (F) {\textbf{action rule} $a = F(c_2, c_3)$\\[-1pt]{\scriptsize accept\,/\,defer\,/\,escalate\,/\,request evidence}};
\draw[arr] (n.south) -- (n.south |- F.north);
\draw[arr] (s.south) -- (s.south |- F.north);
\node[font=\scriptsize, text=black!60, align=center, text width=5.2cm, below=1.5mm of F] (models) {Jev 1.13 via the TypeSafe API, or a comparison model behind the same interface: \mbox{OpenJev-27B}, \mbox{Qwen2.5-7B}, \mbox{OLMo-2-7B}, \mbox{Phi-3.5-mini}};
\draw[arr] (ker.east) -- node[above, font=\scriptsize, pos=0.6]{evidence $x$} (c);

% ---------------- (3) evaluation ----------------
\node[box=red, text width=2.6cm, anchor=west] (y) at ($(F.east)+(11mm,0)$) {\textbf{failure label}\\[-1pt]$y = \mathbf{1}[a \neq a^{*}]$};
\node[box=gray, text width=2.6cm, anchor=south] (ref) at ($(y.north |- s.north)+(0,0)$) {\textbf{target}\\[-1pt]$a^{*} = F(c_2^{*}, c_3^{*})$};
\draw[arr] (ref) -- (y);
\draw[arr] (F.east) ++(0,0) -- node[above, font=\scriptsize, pos=0.55]{action $a$} (y.west);
\node[box=teal, text width=2.6cm, below=3.5mm of y] (rq12) {\textbf{RQ1, RQ2: ranking}\\[-1pt]{\scriptsize AUROC of final-step confidence, action-rule score $r(a)$ and estimator $E_{\mathrm{dec}}$; label-shift reweighting}};
\node[box=violet, text width=2.6cm, below=3mm of rq12] (rq3) {\textbf{RQ3: review value}\\[-1pt]{\scriptsize 10\% review budget, frozen cost matrix; cost vs.\ no review and oracle}};
\draw[arr] (y) -- (rq12);
\draw[arr] (rq12) -- (rq3);

% ---------------- lanes ----------------
\begin{scope}[on background layer]
\node[lane=gray, fit=(moz)(k8s)(cleannote)] (L1) {};
\coordinate (ctop) at ($(c.north)+(0,4.5mm)$);
\node[lane=blue, fit=(c)(s)(F)(models)(ctop)] (L2) {};
\node[lane=teal, fit=(ref)(y)(rq12)(rq3)] (L3) {};
\end{scope}
\node[step=gray, anchor=south west] (t1) at ($(L1.north west)+(1mm,1mm)$) {1}; \node[steplabel=gray] at (t1.east) {Data};
\node[step=blue, anchor=south west] (t2) at ($(L2.north west)+(1mm,1mm)$) {2}; \node[steplabel=blue] at (t2.east) {Typed decision workflow};
\node[step=teal, anchor=south west] (t3) at ($(L3.north west)+(1mm,1mm)$) {3}; \node[steplabel=teal] at (t3.east) {Evaluation};

% reference labels bypass the model
\draw[refarr] (L1.north) -- ++(0,8.5mm) -| node[pos=0.25, above, font=\scriptsize]{reference labels $(c_1^{*}, c_2^{*}, c_3^{*})$, never shown to the model} ($(ref.north)+(8mm,0)$);
\end{tikzpicture}
}
\caption{Study architecture. Reports from three public trackers are
cleaned and stripped of label leaks, then pass through three typed
decisions, each called three times and reduced by modal vote. An action
rule $F$ maps the second and third decisions to an action. Reference
labels, never shown to the model, give the target $\astar$ through the
same rule, and a failure is $a \neq \astar$.}
\label{fig:arch}
\end{widefigure}

\para{Data and workflow} \Cref{fig:arch} shows the pipeline. We use 202
Mozilla and 240 Linux kernel networking and storage bugs for development,
500 newer bugs (250 per tracker, natural label mix) as a held-out set,
and 1{,}082 Kubernetes issues for the replication. Texts are
privacy-scrubbed and truncated to 2{,}000 characters. Node~1 (\choice{})
routes a report to a subsystem, node~2 (\noul{}) asks whether it is a
regression, and node~3 (\score{}) rates severity on three levels. Each
node sees the evidence plus the earlier decisions. Every node is called
three times, and the answers are reduced to one decision by modal vote
(for \noul{}, the mean probability is thresholded at 0.5). All Jev calls
resolved to \texttt{jev-1.13.0}.

\para{Action, target and failure} The rule $F$ maps an unclear $c_2$ to
\emph{request evidence}, $c_3 = 2$ to \emph{escalate}, $c_3 = 1$ with
$c_2 = 1$ to \emph{defer}, and anything else to \emph{accept}. The target
is the same rule on the reference labels, $\astar = F(c_2^{*}, c_3^{*})$,
and a failure is $y = \mathbf{1}[a \neq \astar]$. $F$ is our construct,
so every failure result is conditional on it.

\para{Failure signals} We compare three kinds of signal: the model's
final-step confidence, alone or composed across nodes; the action-rule
score $r(a) = (e(a) + \bar r)/(n(a) + 1)$, the smoothed error rate of the
chosen action on development data; and a logistic estimator on the
decisions ($E_{\text{dec}}$), fitted on development data and frozen.

\para{Evaluation} AUROCs use the Mann--Whitney estimator with 95\%
bootstrap intervals over bugs ($B = 2{,}000$). Confirmatory tests are
one-sided paired bootstrap tests with Holm correction. Costs come from a
matrix frozen in advance: a wrong \emph{accept} costs 6--10, any other
mismatch 1, and each escalation 1 more for the reviewer's time.

\para{Novelty} \Cref{tab:novelty} lists what is new in the design and
what each element lets us conclude.

\begin{coltable}
\centering
\caption{Methodological elements and what each allows us to conclude.}
\label{tab:novelty}
\footnotesize
\begin{tabularx}{\linewidth}{>{\raggedright\arraybackslash}p{2.2cm} Y}
\toprule
Element & What it allows \\
\midrule
Workflow-level failure & Failure is defined on the action taken, so every step-level signal is tested as a predictor of it \\ \addlinespace
Ranking vs.\ cost & A fixed 10\% review budget, a frozen cost matrix and an oracle bound separate ``ranks failures'' from ``reduces cost'' \\ \addlinespace
Shift by design & Development and deployment label mixes differ, which exposes the label-prior dependence of error rates \\ \addlinespace
Pre-registration & Hash-frozen estimators and dated amendments separate confirmatory from exploratory claims \\ \addlinespace
Replication & Five models behind one interface and an independent Kubernetes workflow \\
\bottomrule
\end{tabularx}
\end{coltable}

%======================================================================
\section{Exploratory findings on the development bugs}
\label{sec:explore}
%======================================================================

The three primitives' confidences behave very differently. \choice{}
confidence is near 1.0 for most bugs (85\% of Mozilla and 96\% of kernel
bugs are at $\geq 0.90$), yet it still discriminates its own errors
(AUROC 0.90 on Mozilla, 0.98 on the kernel). \noul{}'s confidence-like
value carries no usable information (AUROC 0.53 and 0.41), and \score{}'s
confidence is about the same at every true level (AUROC 0.58 and 0.57).
Fixed thresholds tuned on constructed data did not transfer: on real
text a \choice{} trigger at 0.90 fired on 85\% and 96\% of bugs, and a
product cut escalated 88\% and 78\% against targets of 41\% and 42\%.
\begin{coltable}
\centering
\caption{Node agreement with reference labels on the development bugs
(Wilson 95\% intervals). Baseline: the most frequent label.}
\label{tab:agree}
\footnotesize
\setlength{\tabcolsep}{3.5pt}
\begin{tabular}{lcccc}
\toprule
 & \multicolumn{2}{c}{Mozilla} & \multicolumn{2}{c}{Kernel} \\
\cmidrule(lr){2-3}\cmidrule(lr){4-5}
Node & Agreement & Base & Agreement & Base \\
\midrule
\choice{} & 0.89 [0.83, 0.92] & 0.59 & 0.975 [0.95, 0.99] & 0.50 \\
\noul{}   & 0.80 [0.74, 0.85] & 0.74 & 0.725 [0.67, 0.78] & 0.81 \\
\score{}  & 0.53 [0.46, 0.60] & 0.40 & 0.44 [0.38, 0.50] & 0.33 \\
\bottomrule
\end{tabular}
\end{coltable}

The answers themselves are uneven (\cref{tab:agree}). \choice{} routes
well above baseline on both trackers. \noul{}'s errors change direction:
on Mozilla it mostly misses regressions (23 misses vs.\ 17 false
alarms), while on the kernel it over-flags them (63 false alarms vs.\ 3
misses). \score{} compresses severity toward the middle level
(\cref{fig:sev}). Most minor bugs are rated higher than their label
(79\% on Mozilla, 75\% on the kernel), and many serious bugs lower (43\%
and 56\%). Using Mozilla's own severity definitions moved agreement only
from 0.53 to 0.56, so the compression does not come from our mapping.

\begin{colfigure}
\centering
\begin{tikzpicture}
\begin{groupplot}[group style={group size=2 by 1, horizontal sep=0.35cm, y descriptions at=edge left},
  width=0.5\linewidth, height=3.8cm, view={0}{90}, enlargelimits=false, axis on top,
  colormap={sev}{rgb255=(255,255,255) rgb255=(107,174,214) rgb255=(33,113,181)},
  point meta min=0, point meta max=0.8,
  xtick={0,1,2}, ytick={0,1,2}, y dir=reverse,
  yticklabels={minor,normal,serious}, xticklabels={0,1,2},
  tick label style={font=\scriptsize}, label style={font=\scriptsize}, title style={font=\small},
  xlabel={\score{} output},
  nodes near coords={\pgfmathprintnumber[fixed,fixed zerofill,precision=2]{\pgfplotspointmeta}},
  nodes near coords style={font=\scriptsize, anchor=center}]
\nextgroupplot[title={Mozilla (triage)}, ylabel={true level}]
\addplot[matrix plot*, mesh/cols=3, point meta=explicit] table[meta=v] {fig_heat_mozilla.dat};
\nextgroupplot[title={Kernel (reporter)}]
\addplot[matrix plot*, mesh/cols=3, point meta=explicit] table[meta=v] {fig_heat_kernel.dat};
\end{groupplot}
\end{tikzpicture}
\caption{Row-normalized confusion of \score{} against reference
severity. On both trackers the middle column absorbs mass from both
extremes.}
\label{fig:sev}
\end{colfigure}

Composing confidences did not help. A naive comparison looks
favourable to decisions: on Mozilla, the decision-based action beat a
confidence-product rule by 0.17 [0.08, 0.27] in the probability of the
right action (exact McNemar $p = 0.0004$). But the confidence rule never
reads the decisions, while the target escalates on the severity
decision, so that comparison is confounded. The fairer test starts from
the decision-based action and adds escalations ranked by low composed
confidence. Among
bugs the rule did not escalate, composed confidence ranked missed serious
bugs at chance (AUROC 0.49 and 0.50), and across all bugs it ranked wrong
actions at chance as well (0.54 and 0.47). At a budget of 10\% of bugs,
confidence-ranked overrides did no better than random selection on
either tracker. Every one of these intervals
covers 0.5 (\cref{fig:forest}). The reason is structural. The only
informative confidence belongs to \choice{}, the one node $F$ does not
read.

\begin{colfigure}
\centering
\begin{tikzpicture}
\begin{axis}[width=0.78\linewidth, height=4.6cm, xmin=0.25, xmax=1.05, ymin=0.4, ymax=5.6,
  ytick={1,2,3,4,5}, yticklabels={(B) any error, (A) missed serious, \score{} own, \noul{} own, \choice{} own},
  xlabel={AUROC (95\% bootstrap interval)}, grid=major, grid style={gray!20},
  tick label style={font=\scriptsize}, label style={font=\scriptsize},
  legend style={font=\scriptsize, at={(0.98,0.02)}, anchor=south east}]
\draw[dashed, gray] (axis cs:0.5,0.4) -- (axis cs:0.5,5.6);
\addplot[blue, only marks, mark=*, error bars/.cd, x dir=both, x explicit]
  table[x=x, y=y, x error minus=lo, x error plus=hi] {fig_forest_mozilla.dat};
\addplot[orange, only marks, mark=square*, error bars/.cd, x dir=both, x explicit]
  table[x=x, y=y, x error minus=lo, x error plus=hi] {fig_forest_kernel.dat};
\legend{Mozilla, Kernel}
\end{axis}
\end{tikzpicture}
\caption{Confidence AUROCs on the development bugs; dashed line: chance.
Only \choice{}'s confidence discriminates its own errors, and composed
confidence is at chance for both escalation targets, (A) missed serious
bugs and (B) any wrong action.}
\label{fig:forest}
\end{colfigure} These observations led to the pre-registered tests below.

%======================================================================
\section{Pre-registered held-out test}
\label{sec:heldout}
%======================================================================

We trained $E_{\text{dec}}$ ($L_2$-regularized logistic regression on
the \score{} level, the \noul{} decision and the action) on the 442
development bugs, froze it by hash, and only then collected the 500
held-out bugs (4{,}500 Jev calls). We pre-registered two hypotheses:
$E_{\text{dec}}$ ranks wrong actions better than low final-step
confidence (HR1), and above chance on each tracker (HR2).

Both were supported. $E_{\text{dec}}$ reached AUROC 0.763 against 0.549
for confidence, a difference of $+0.214$ with a one-sided lower bound of
$+0.158$. Per tracker it reached 0.811 on Mozilla and 0.697 on the
kernel, and the gap over confidence was $+0.279$ [$+0.191$, $+0.374$] on
Mozilla and a weaker $+0.124$ [$+0.023$, $+0.225$] on the kernel. On the
70 Mozilla bugs filed in the 12 months before sampling, which are least
likely to appear in model training data, $E_{\text{dec}}$ reached 0.847
against 0.519. The result held when every development bug linked to a held-out
bug was removed and the estimator refitted (0.772 against 0.551), and an
independent script that shares no code with the analysis reproduced
every point estimate. The estimator's predicted failure rate matches
the observed rate overall, but it outputs only 7--8 distinct risk levels
(ECE 0.18--0.20), so it ranks better than it calibrates.

%======================================================================
\section{Cross-model replication}
\label{sec:crossmodel}
%======================================================================

To see whether the result is specific to Jev, we ran the same bugs,
questions and adapter through four more models: OpenJev-27B
\citep{openjev2026}, an independently trained Jev-style model, and three
small general instruct models, Qwen2.5-7B \citep{qwen2024qwen25},
OLMo-2-7B \citep{olmo2024olmo2} and Phi-3.5-mini \citep{abdin2024phi3}.
Jev is used zero-shot, so the fair comparison is other models used the
same way, through the same questions and TypeSafe's confidence formulas.
We picked small open models because they are cheap, reproducible and
make more mistakes, which gives the failure predictors more to find. The
protocol had three amendments: OpenJev was added before any of its
outputs existed, OLMo-2 needed a newer \texttt{transformers} release,
and a readout bug found after interim results led us to re-run all three
general models. The conclusions did not change.

\begin{colfigure}
\centering
\begin{tikzpicture}
\begin{axis}[width=0.8\linewidth, height=4.2cm, xmin=0.3, xmax=0.9, ymin=0.4, ymax=5.6,
  ytick={1,2,3,4,5}, yticklabels={OLMo-2, Phi-3.5, Qwen2.5, OpenJev, Jev 1.13},
  xlabel={held-out AUROC for wrong actions (95\% CI)}, grid=major, grid style={gray!20},
  tick label style={font=\scriptsize}, label style={font=\scriptsize},
  legend style={font=\scriptsize, at={(0.02,0.02)}, anchor=south west}]
\draw[dashed, gray] (axis cs:0.5,0.4) -- (axis cs:0.5,5.6);
\addplot[blue, only marks, mark=*, error bars/.cd, x dir=both, x explicit]
  table[x=dec, y expr=\thisrow{y}+0.12, x error minus=declo, x error plus=dechi] {fig_crossmodel.dat};
\addplot[orange, only marks, mark=square*, error bars/.cd, x dir=both, x explicit]
  table[x=conf, y expr=\thisrow{y}-0.12, x error minus=conflo, x error plus=confhi] {fig_crossmodel.dat};
\legend{decision estimator, final-step confidence}
\end{axis}
\end{tikzpicture}
\caption{For all five models, the decision estimator ranks wrong
workflow actions better than the model's final-step confidence on the
500 held-out bugs.}
\label{fig:crossmodel}
\end{colfigure}

\begin{widetable}
\centering
\caption{Per-model results on the 500 held-out bugs. AUROC columns rank
wrong workflow actions (95\% bootstrap CI): the decision estimator fitted
on each model's own development outputs, final-step confidence, the
action-rule score and Jev's frozen estimator applied unchanged. Cost
columns are from the pre-registered review test (\cref{sec:intervention});
the oracle is an upper bound that must spend all 50 reviews.}
\label{tab:crossmodel}
\resizebox{\textwidth}{!}{% Generated by make_figure_data.py from research/results/*.json. Do not edit by hand.
\begin{tabular}{llcccccrrrr}
\toprule
 & & & \multicolumn{4}{c}{Held-out AUROC for wrong actions} & \multicolumn{4}{c}{Cost at 10\% review budget} \\
\cmidrule(lr){4-7}\cmidrule(lr){8-11}
Model & Type & Fail rate & Decision est. & Final-step conf. & Action rule & Jev est. & No review & Best review arm & Oracle & Wrong accepts \\
\midrule
Jev & TypeSafe, hosted & 0.43 & 0.76 [0.72, 0.80] & 0.55 [0.49, 0.60] & 0.76 & --- & 658 & 681 (conf.\ product) & 370 & 52\% \\
OpenJev & Jev-style, 27B & 0.47 & 0.75 [0.70, 0.79] & 0.57 [0.53, 0.63] & 0.75 & 0.76 & 660 & 674 (confidence) & 441 & 40\% \\
Qwen2.5 & general, 7B & 0.74 & 0.74 [0.68, 0.79] & 0.47 [0.41, 0.54] & 0.70 & 0.74 & 903 & 910 (estimator) & 856 & 14\% \\
Phi-3.5 & general, 3.8B & 0.72 & 0.68 [0.62, 0.73] & 0.40 [0.34, 0.46] & 0.69 & 0.70 & 831 & 890 (confidence) & 890 & 4\% \\
OLMo-2 & general, 7B & 0.80 & 0.64 [0.58, 0.69] & 0.55 [0.48, 0.62] & 0.58 & 0.64 & 893 & 916 (all arms tie) & 916 & 1\% \\
\bottomrule
\end{tabular}
}
\end{widetable}

All 8 pre-registered tests were supported (\cref{fig:crossmodel},
\cref{tab:crossmodel}).
Decision estimators reached AUROC 0.64--0.76, final-step confidence
0.40--0.58. The weakest case was OLMo-2 (0.637 vs.\ 0.551), which passed
at the last Holm step. Jev's frozen estimator, never refitted, reached
0.76 [0.72, 0.80] on OpenJev's outputs, as well as OpenJev's own
estimator did.

Most of this signal is simply which action was chosen. The action's
error rate on development data matches the estimator for Jev (0.76) and
OpenJev (0.75) and comes close for the others (0.58--0.70). For Jev the
estimator adds $+0.000$ [$-0.012$, $+0.011$] beyond that rate. For Qwen2.5 and OLMo-2 it adds a little
($+0.035$ and $+0.057$). Single uncertainty signals sit at chance:
routing confidence at 0.48--0.52, regression confidence at 0.45--0.56
and severity entropy at 0.40--0.55. Adding them to the decision features
helps only the weakest general models ($+0.054$ for Phi-3.5, $+0.078$ for
OLMo-2). Part of this is built in, since the
action and the target come from the same rule $F$. The two model
families also fail differently. Jev and OpenJev pull severity to the
middle level, while the general models call 67--100\% of bugs
``blocking'' whatever the true level, and are wrong far more often
(72--80\% vs.\ 43--47\%).

That makes the rates sensitive to the label mix (RQ2). The development
set was quota-sampled by severity and the held-out set was not. On a
target that bypasses $F$ (the severity answer being wrong), estimators
fitted on development data ranked held-out errors far below chance, 0.28
for Jev and OpenJev. Jev's level-1 answers were 60\% wrong in the
development set and 20\% wrong in the natural mix. Re-weighting to the
natural mix fixed this for four of five models (\cref{tab:targets}; Jev
0.79 vs.\ 0.61 for confidence).

\begin{coltable}
\centering
\caption{Exploratory: held-out AUROC of the decision estimator on two
targets that bypass $F$, fitted on the quota-sampled development bugs
(``dev fit'') or re-weighted to the natural label mix, against
final-step confidence. Bold: re-weighted estimator above confidence.}
\label{tab:targets}
\small
\resizebox{\linewidth}{!}{% Generated by make_figure_data.py from baseline_ablations.json and prior_shift_reweighting.json. Do not edit by hand.
\begin{tabular}{lcccccc}
\toprule
 & \multicolumn{3}{c}{Severity wrong} & \multicolumn{3}{c}{Any node wrong} \\
\cmidrule(lr){2-4}\cmidrule(lr){5-7}
Model & dev fit & re-weighted & conf. & dev fit & re-weighted & conf. \\
\midrule
Jev & 0.28 & \textbf{0.79} & 0.61 & 0.56 & \textbf{0.76} & 0.56 \\
OpenJev & 0.28 & \textbf{0.79} & 0.63 & 0.62 & \textbf{0.81} & 0.56 \\
Qwen2.5 & 0.62 & \textbf{0.74} & 0.46 & 0.76 & \textbf{0.73} & 0.40 \\
Phi-3.5 & 0.31 & \textbf{0.67} & 0.41 & 0.67 & \textbf{0.77} & 0.29 \\
OLMo-2 & 0.50 & \textbf{0.62} & 0.57 & 0.42 & 0.42 & 0.49 \\
\bottomrule
\end{tabular}
}
\end{coltable}
 Estimating
the mix without labels by BBSE \citep{lipton2018bbse} recovered most of
that gain for Jev (0.74), OpenJev (0.78) and Qwen2.5 (0.74), but not for
the right reason. The estimated mix was wrong. For the kernel it was
$[0, 1, 0]$ against a true $[0.02, 0.73, 0.25]$, and the ranking improved
only because the middle level happened to dominate. Under known
synthetic shifts the estimate was biased toward the development mix, and
for Phi-3.5, whose severity answers are nearly all ``blocking'', it gave
no gain at all.

%======================================================================
\section{An independent workflow: Kubernetes}
\label{sec:k8s}
%======================================================================

We pre-registered a replication on issue triage in the Kubernetes GitHub
repository before collecting any of its data. It differs in platform,
labelling process, routing (network, storage, node) and questions (bug
or not, priority). The split is chronological: 500 issues from 2015--2019
for development and the 582 latest eligible issues from 2021--2026 for
testing, with linked issues removed and label leaks stripped from the
text.

\begin{coltable}
\centering
\caption{Kubernetes test set: AUROC for wrong actions of the action-rule
score and of final-step confidence, and priority accuracy (chance 0.33).}
\label{tab:k8s}
\footnotesize
\begin{tabular}{lcccl}
\toprule
Model & Rule & Conf. & Priority acc. & Result \\
\midrule
Jev 1.13     & 0.776 & 0.628 & 0.371 & supported \\
OpenJev-27B  & 0.796 & 0.569 & 0.404 & supported \\
Qwen2.5-7B   & 0.710 & 0.555 & 0.270 & supported \\
OLMo-2-7B    & 0.850 & 0.255 & 0.082 & supported \\
Phi-3.5-mini & 0.519 & 0.674 & 0.046 & not replicated \\
\bottomrule
\end{tabular}
\end{coltable}

For Jev the action rule again beat confidence, 0.776 against 0.628, a
paired difference of $+0.148$ [$+0.091$, $+0.207$]. As before, the
decision estimator added little ($+0.004$ to $+0.017$), and the contrast
held for the any-node target, within SIG Node, and on issues triaged
within 30 days. Because the two periods also differ in issue format and
label mix, this is a combined test of temporal and distribution shift. A second
pre-registration ran the four other models (\cref{tab:k8s}). By its fixed
rule, at least three of four, the finding holds. The detail matters,
though. Only Jev and OpenJev-27B beat chance on every step. The three
general models handled routing and the bug question but failed on
priority, so their workflows were wrong most of the time (wrong-action
rates up to 95\%), and the rule ``won'' largely by flagging the one
action each model nearly always chose. Phi-3.5 escalated 580 of the 582
test issues, which left the rule nothing to rank.

The label-free correction did not carry over. Re-weighting a deliberately
mismatched fit by BBSE lowered its AUROC from 0.699 to 0.348. BBSE again
estimated the mix as entirely middle level, $[0, 1, 0]$, while the true
mix was $[0.69, 0.26, 0.05]$, so the weighting pushed the fit the wrong
way. This supports our reading of the earlier gains as a coincidence of
a dominant middle level.

%======================================================================
\section{Does better ranking reduce cost?}
\label{sec:intervention}
%======================================================================

We pre-registered the cost test before computing any cost outcome. Each
model gets 50 reviews (10\% of the held-out bugs), and a review sends the
case to a person. The primary arm ranks candidates by the expected
saving of a review for their action, estimated on development data. It
is compared with confidence-ranked and random review.

\begin{colfigure}
\centering
\begin{tikzpicture}
\begin{axis}[ybar, bar width=3pt, width=\linewidth, height=4.0cm,
  symbolic x coords={Jev,OpenJev,Qwen2.5,Phi-3.5,OLMo-2}, xtick=data, ymin=0, ymax=960,
  ylabel={total cost (500 bugs)}, enlarge x limits=0.12,
  grid=major, grid style={gray!20}, tick label style={font=\scriptsize}, label style={font=\scriptsize},
  legend image code/.code={\draw[#1] (0cm,-0.08cm) rectangle (0.18cm,0.12cm);}, legend style={font=\scriptsize, at={(0.5,1.02)}, anchor=south, legend columns=5, draw=none}]
\addplot[fill=gray!40] table[x=model, y=noreview] {fig_intervention.dat};
\addplot[fill=blue!55] table[x=model, y=A1] {fig_intervention.dat};
\addplot[fill=orange!70] table[x=model, y=B] {fig_intervention.dat};
\addplot[fill=gray!75] table[x=model, y=C] {fig_intervention.dat};
\addplot[fill=green!55!black] table[x=model, y=oracle] {fig_intervention.dat};
\legend{no review, action rule, confidence, random, oracle}
\end{axis}
\end{tikzpicture}
\caption{Total cost at a 10\% review budget under the pre-registered cost
matrix. No feasible strategy beats no review. A perfect reviewer (oracle)
would cut Jev's cost by 44\% and OpenJev's by 33\%.}
\label{fig:intervention}
\end{colfigure}

Two of ten pre-registered tests were supported: the action rule beat
random review for Jev and OpenJev, by about 10--12 cost units. It never
beat confidence, and every review arm cost more than no review at all.
For Jev that was 658 with no review against 681--714 with review
(\cref{fig:intervention}).

The oracle shows where the money is. For the Jev-style models, 40--52\%
of no-review cost comes from wrongly accepted serious or regression
bugs, and picking the right 50 cases would save 33--44\%. Those cases are
only 15--16\% of accepted cases, and nothing we tested could find them.
Among Jev's 245 accepted held-out cases, the decision estimator,
final-step confidence and \score{}'s top-level probability ranked the 38
wrong accepts at AUROC 0.46, 0.49 and 0.49. For the general models the
problem is different: they over-escalate, and even the oracle saves at
most 5\%.

\begin{coltable}
\centering
\caption{Exploratory cost sensitivity: best feasible arm minus no review
(negative, in bold: review reduces cost), rankings frozen. $m$ scales
wrong-accept costs, $h$ is the human cost per escalation and $p^{*}$ the
break-even precision. $^{\dagger}$The pre-registered table.}
\label{tab:costsens}
\small
\resizebox{\linewidth}{!}{% Generated by make_figure_data.py from research/results/intervention_cost_sensitivity.json. Do not edit by hand.
\begin{tabular}{rrrrrrrr}
\toprule
$m$ & $h$ & $p^{*}$ & Jev & OpenJev & Qwen2.5 & Phi-3.5 & OLMo-2 \\
\midrule
0.5 & 0.5 & 25\% & +29 & +25 & +21 & +52 & +20 \\
0.5 & 1 & 33\% & +54 & +50 & +46 & +77 & +28 \\
0.5 & 2 & 50\% & +104 & +100 & +96 & +127 & +45 \\
1 & 0.5 & 14\% & \textbf{-2} & \textbf{-11} & \textbf{-18} & +34 & +14 \\
1 & 1$^{\dagger}$ & 18\% & +23 & +14 & +7 & +59 & +23 \\
1 & 2 & 27\% & +73 & +64 & +57 & +109 & +40 \\
2 & 0.5 & 7\% & \textbf{-69} & \textbf{-83} & \textbf{-97} & \textbf{-2} & +4 \\
2 & 1 & 10\% & \textbf{-44} & \textbf{-58} & \textbf{-72} & +23 & +13 \\
2 & 2 & 14\% & +6 & \textbf{-8} & \textbf{-22} & +73 & +30 \\
\bottomrule
\end{tabular}
}
\end{coltable}

The conclusion depends on the costs (\cref{tab:costsens}). In an
exploratory re-scoring within a factor of 2 of the frozen matrix, review
pays for Jev, OpenJev and Qwen2.5 when a wrong accept costs twice as much
(for Jev, 44--69 units below no review) or a review half as much,
because the break-even precision then falls to about the 14--16\% share
of wrong accepts. The ordering of the arms is steadier than their sign:
the action rule beats random for both Jev-style models in 12 of 16
cells. If reviewers are right only 80\% of the time, Jev's saving at twice
the wrong-accept cost falls from 44 to 26 units but does not
disappear.
We also tried a per-case score built to find the expensive errors. On
the development mix it cut cost by 12\% against no review (630 vs.\ 713),
but the development bugs were quota-sampled, and re-weighted to the
natural mix it was projected to cost more than no review. Its
pre-registered feasibility check required power of at least 0.80 with at
most 1{,}000 fresh bugs. The highest power was 0.20, so we stopped that
study before collecting new data.

%======================================================================
\section{Lessons from our own audits}
\label{sec:method}
%======================================================================

Our audit process caught seven defects, five of them before the data
that would have exposed them existed. Among them were a
``primitive-aware'' strategy that turned out to be identical to the
baseline, a trigger that fired on about 90\% of inputs, label leaks in
constructed content, and a bootstrap seed that made two reported bounds
vary between runs by up to 0.005. The confidence-versus-decisions
confound was found only after a result. It would have been caught
earlier by asking, for each strategy and each target class, whether the
strategy had the information it needed to be right. We suggest that
check, together with measuring each trigger's selectivity on pilot
inputs, before any escalation policy is frozen.

%======================================================================
\section{Threats to validity}
\label{sec:threats}
%======================================================================

The largest threat is construct validity. Failure means disagreeing with
a target built from tracker labels, and those labels are not ground
truth: Mozilla's sample excludes re-triaged bugs, kernel severities are
reporters' claims, and the kernel's middle level is a default value. In an exploratory pilot, one AI assistant chose actions for 50
held-out bugs and agreed with the target at about chance (Cohen's
$\kappa$ between $-0.157$ and $+0.015$). With one rater we cannot tell
whether the target or the judge is at fault. A two-human validation is
designed but not yet run. The cost matrix is also our construct, and real
review costs are unknown.

On internal validity, the cost test was pre-registered after the
prediction results on the same held-out bugs were known, though before
any cost was computed, and the cost-sensitivity and label-shift analyses
are exploratory. Old public bugs may appear in model
training data, and the recent subsets we can check are small (Mozilla 54,
kernel 35). The
protocols were revised several times, including a readout bug in the
cross-model study fixed after interim results were seen. Every revision
is dated and disclosed.

On external validity, the study covers five models in one domain,
software issue triage, across two workflows and one Jev version. On
Kubernetes the general models were below chance on priority, and the
Kubernetes leakage read was done by an AI assistant at our direction.

%======================================================================
\section{Conclusion}
\label{sec:conclusion}
%======================================================================

In a chain of three typed decisions, what the workflow decided predicted
its failures far better than how confident it was. This held on 500 new
bugs, for five models, and on an independent Kubernetes workflow. Most
of the signal is the error rate of the chosen action, and that rate
depends on the label mix it was measured under. Predicting failures did
not make review worthwhile. The costly errors are wrong accepts that
look like correct accepts to every signal we tested, so whether review
pays comes down to the ratio of error cost to review cost and to how
precisely a signal can find those cases.

%======================================================================
\section{Future scope}
\label{sec:future}
%======================================================================

The first step is to validate the target with two independent human
evaluators. Beyond that, the oracle points to the real opportunity: a
signal that separates costly errors within a single action. OpenJev's
weak within-action severity signal (AUROC 0.68) is a place to start.
The study should also be repeated outside software triage, with newer
model versions, supervised baselines and measured review costs.
